\subsection{Supervised Data and Classifier Controls}
\label{sec:supervised_methods}

To assess whether an LLM is warranted when labelled data are available,
we add TF-IDF/LinearSVC and fine-tuned DistilBERT
\cite{sanh2019distilbert} baselines. All twelve saved test sets remain
unchanged. Training and validation use separate balanced sets of
20,000 and 2,000 records, respectively. We exclude normalised exact
matches and shared 250-character prefixes with any test narrative,
discard 133 conflicting-label text groups, and screen near duplicates
using 128-permutation MinHash retrieval followed by exact word-trigram
Jaccard similarity at threshold 0.85. This excludes 496 additional
test-similar records and 28 validation-similar training candidates.
The approximate retrieval cannot guarantee removal of every paraphrase.

The preserved test narratives are reserved before drawing training and
validation records. Test-text matching is an exclusion constraint during
split construction, not a source of fitted features. Verification finds
zero normalised exact or shared 250-character-prefix matches between
training, validation and tests at all four training budgets. Identity checks
use SHA-256 hashes of normalised text and its 250-character prefix;
the saved source index is a local CSV index, not a CFPB complaint identifier.
Split hashes are recorded in a manifest; test-file hashes are rechecked
before paired HF evaluation. Supplementary Section~\ref{S-sec:split_integrity}
details the split counts, sampling seeds and statistical units.

Nested balanced training subsets of 500, 2,000, 10,000, and 20,000
records are used for the classical learning curve. TF-IDF uses word
unigrams/bigrams, sublinear term frequencies, minimum document
frequency 2, and at most 100,000 features. Vocabulary and IDF weights
are fitted exclusively on each training subset; validation and test
inputs are transformed without refitting. This differs from the fixed,
non-learned lowercasing and character/whitespace normalisation.
LinearSVC's $C$ is selected
from $\{0.1,1,10\}$ by validation weighted F1. Warm serial CPU latency includes vectorisation and prediction.

The encoder is \texttt{distilbert/distilbert-base-uncased}, pinned to
the revision listed in the supplementary material. The 20,000-record
and fixed 500-record controls train for three epochs with learning rate
$2\times10^{-5}$, batch size 16, weight decay 0.01, a 256-token cap,
and 10\,\% warmup.
The checkpoint with highest validation weighted F1 is selected.
Training seeds are distinct from the three saved test-subset identifiers.
Within each supervised fit, checkpoint and hyperparameter selection
use the separate validation set, not test performance.
The training and prediction artifacts report exact model commits,
training time, peak allocated VRAM, and warm serial GPU latency.

We additionally train DistilBERT on the identical 500-record subset
with three training seeds and the same 2,000-record validation set.
A fixed three-epoch control had a validation trajectory that was still
improving, motivating a separate
convergence control before inspecting its test scores. This control
allows at most 30 epochs, uses patience 5 on validation weighted F1,
and selects the best validation checkpoint. All other encoder settings
are retained; original three-epoch outputs are preserved separately.
This matches label budgets, not training-update or compute budgets.

All supervised budgets use an additional 2,000 validation labels: total
label counts are 2,500, 4,000, 12,000 and 22,000 for the classical
training budgets. Each LLM adapter fit likewise uses 2,500 labels in
total. A training budget is not the total annotation requirement.
Matching label availability does not equate validation objectives,
fit counts or compute effort (Supplementary Section~\ref{S-sec:split_integrity}).
